\section{Introduction}\label{sec:introduction}

Modern IoT and enterprise systems generate large volumes of interactions among devices, hosts, services, processes, and files. Intrusion detection in these environments is difficult because attacks are often rare, evolving, and only partially represented in labeled datasets. Traditional rule-based intrusion detection systems depend on known signatures, while conventional statistical detectors often rely on fixed assumptions about normal and abnormal behavior. These approaches can struggle when facing zero-day attacks, where the malicious behavior has not been observed before and does not match predefined rules or known attack patterns \cite{Buczak2016,sommer2010outside}. Machine learning-based intrusion detection is therefore attractive because it can learn behavioral representations from data and identify suspicious deviations even when the exact attack type is unknown \cite{Buczak2016}.

However, fully supervised machine learning also has an important limitation for intrusion detection: it usually requires enough labeled examples of both non-malicious and malicious behavior. In practice, confirmed malicious samples are rare, expensive to label, and often become outdated as attacker behavior changes \cite{sommer2010outside,caville2022anomale}. A model trained only to recognize known attacks may fail when a new attack does not resemble its labeled training examples \cite{sommer2010outside}. This motivates an unsupervised, or benign-only, learning setting. Instead of learning attack-specific decision boundaries, the detector learns a generalized representation of normal system behavior and flags graphs that deviate from the distribution as potential intrusions \cite{manzoor2016streamspot,caville2022anomale}. This setting is especially useful for zero-day detection because it does not require attack labels during representation learning.

System activity is inherently relational, which makes graphs a natural representation for intrusion detection. A single event or flow may appear harmless when inspected alone, but its role in a larger interaction pattern can reveal malicious behavior. Port scanning can generate one-to-many connection patterns \cite{lee2003detection}. Lateral movement can create multi-step paths across hosts rather than isolated events \cite{fang2022lmtracker,zhou2024lmdetect}. Denial-of-service and coordinated attacks can concentrate repeated requests or flows around selected targets \cite{bakar2024ftgnet}. In graph form, these behaviors may appear as star-like subgraphs, branches, paths, or other structural changes across hosts and services \cite{iliofotou2007tdg}. Graph-based intrusion detection models these interactions explicitly: nodes represent entities such as hosts, devices, services, processes, or files, and edges represent interactions among them. A graph-based detector can therefore judge an event through its surrounding structure instead of treating each record as independent \cite{manzoor2016streamspot,lo2022egraphsage}.

Topology is important for graph-based intrusion detection for two reasons. First, it supports provenance and explainability. Many attacks unfold as paths through a system: an attacker compromises one entity, uses it to reach another, and gradually expands control. Advanced persistent threat activity and lateral movement can therefore leave structural traces in the form of paths, branches, and connectivity changes across processes, files, services, and hosts \cite{fang2022lmtracker,zhou2024lmdetect}. These topological signatures can help explain why a graph is suspicious by relating the anomaly to the shape of the interaction pattern. Second, topology captures global structure. Intrusions may not always appear as obviously abnormal individual events; instead, they may change how parts of the system are connected. A detector that only focuses on local neighborhoods can miss these broader behavioral relationships, while topology-aware representations can summarize graph shape across multiple scales.

Graph contrastive learning provides a useful way to learn such representations without attack labels. It creates augmented views of the same graph and trains an encoder to recognize that the views come from the same original graph. GraphCL showed that this approach can learn useful graph-level representations without supervised labels \cite{you2020graphcl}. However, standard GNN message passing mainly aggregates local neighborhoods. In intrusion detection, this can be limiting because non-malicious and malicious graphs may share the same broad hub-and-spoke structure while differing in a few security-critical paths, branches, or connections. TopoGCL addresses this problem by adding topological summaries that describe graph shape across multiple scales, making topology-aware contrastive learning a strong starting point for graph intrusion detection \cite{chen2024topogcl}.

Despite this promise, TopoGCL was designed mainly for unsupervised graph classification, not for anomaly scoring in an intrusion detection system \cite{chen2024topogcl}. Classification and intrusion detection have different goals. In graph classification, labeled examples can help the model learn class boundaries after representation learning. In intrusion detection, there are no labeled attack examples during representation learning. The detector must decide whether a new graph is abnormal by comparing it with the learned distribution. Because of this, the embedding space used for distance-based scoring directly controls detection quality.

We propose \sysname, a topology-aware graph contrastive learning method for intrusion detection. \sysname combines unsupervised learning, topological graph representation, and contrastive training to learn embeddings that better separate normal system activity from structurally abnormal behavior. During training, \sysname learns from augmented views of non-malicious graphs only. It combines graph-level message-passing embeddings with topology-aware structural embeddings. During detection, it scores a test graph by measuring its distance from non-malicious training embeddings. Graphs that are farther from normal training graphs receive higher anomaly scores.

The key design choice in \sysname is to score graphs in the same projected embedding space used by contrastive training. Many contrastive learning pipelines train a projection head but later compute anomaly scores from the encoder's raw graph embeddings. This creates a mismatch between the space optimized by the contrastive objective and the space used for detection. \sysname avoids this mismatch by comparing graphs after projection. This change does not require attack labels during representation learning, an extra classifier, or a much larger model.

We evaluate \sysname on StreamSpot, GraSec-IoT, and Unicorn Wget \cite{manzoor2016streamspot,hamouche2024graseciot,unicornwget}. StreamSpot contains heterogeneous system-activity graphs. GraSec-IoT contains communication and attack behavior from an Internet of Things setting. Unicorn Wget contains graph samples from a system-activity attack scenario. Together, these datasets test whether \sysname works across different graph construction settings and attack behaviors. We compare \sysname with a supervised GNN and GraphCL, a graph contrastive learning baseline.

In summary, this paper makes the following contributions:
\begin{itemize}
    \item We study ML-based graph intrusion detection for intelligent IoT and enterprise activity graphs under benign-only representation learning. This setting is important for zero-day detection because malicious samples are rare, labels are costly, and new attacks may not resemble known attack examples.

    \item We propose \sysname, a topology-aware graph contrastive learning method for intrusion detection. \sysname uses topological information to capture provenance-style attack paths and global structural relationships that may be missed by purely local message passing.

    \item We evaluate \sysname on StreamSpot, GraSec-IoT, and Unicorn Wget against a supervised GNN and GraphCL. \sysname improves F1 score from $0.923$ to $0.926$ on StreamSpot, from $0.790$ to $0.985$ on GraSec-IoT, and from $0.583$ to $0.863$ on Unicorn Wget.
\end{itemize}

Overall, \sysname shows that unsupervised topology-aware contrastive learning can provide stronger graph representations for intrusion detection by learning normal system behavior, preserving global structural information, and scoring deviations in the embedding space optimized during training.
